\uuid{XCzK}
\titre{Distance}
\auteur{bodin ; hanani ; mzari}
\organisation{exo7}
\licence{CC BY-NC-SA 4.0}
\datecreate{2020-01-20}
\niveau{L1}
\module{Géométrie}
\chapitre{Géométrie de l'espace}
\sousChapitre{Distance}
\difficulte{3}
\theme{Distance}
\exo7qcmid{lille-1.geometrie-espace.027}

\contenu{
\texte{
Pour ces questions, l'espace est muni d'un repère orthonormé direct  $(O, \vec{i}, \vec{j}, \vec{k})$.
}

\begin{question}
\qclassification{141.05}
\qtags{motcle=distance}
Soit  $A(1,1,1)$ un point,  $D$ la droite  $ : \left\{\begin{array}{ccl}x&=&1+z\\y&=&z \end{array}\right.$ et $P$ un plan contenant $D$ et tel que la distance de $A$ à $P$ soit égale à   $\frac{1}{\sqrt 2}$. Une équation cartésienne de $P$ est :
\begin{answers}

       \bad {$x+z+1=0$ ou $x+y+1=0$}

      \bad {$x-z+1=0$ ou $x-y=0$}

     \bad {$z=1$ ou $x=1$}

      \good {$x-z=1$ ou $x-y=1$}

\end{answers}
\begin{explanations}
$P$ est différent du plan d'équation  $y-z=0$ et $D \subset P$, une équation cartésienne de $P$ est donc de la forme : $(x-z-1)+ \alpha (y-z)=0$. On calcule $\alpha$ pour que la distance de $A$ à $P$ soit égal à $\frac{1}{\sqrt 2}$.
\end{explanations}
\end{question}
}
